% Part: computability
% Chapter: computability-theory
% Section: reducibility

\documentclass[../../../include/open-logic-section]{subfiles}

\begin{document}

\olfileid{cmp}{thy}{red}
\olsection{قابلية الاختزال}

\begin{explain}
قد حصل لنا مثال على الأقل، هو $K_0$، لمجموعة تقبل التعداد الحسابي
ولا تقبل الحساب، وليست منفردة بهذه الصفة. وبطريقة
الاختزال نثبت هذه الخواص لمجموعات أخرى، من غير أن نستأنف
في كل مرة البرهان من أصوله.

ومعنى «اختزال» المجموعة $A$ إلى المجموعة~$B$، على العموم، أن نستخرج من أجوبة
السؤال عن انتماء ال!!{element}s إلى~$B$ أجوبة السؤال عن انتمائها
إلى~$A$. ونخص بالبحث هنا «قابلية الاختزال متعددًا إلى واحد»؛
وللاختزال مفاهيم أخرى كثيرة تختلف خواصها. وللاختزال
منزلة أساسية في التعقيد الحسابي أيضًا، حيث يُنظر، فوق إمكان الحساب، في
كفاءته. فالمجموعة «تامة في NP» إذا كانت في NP وأمكن
اختزال كل مسألة في NP إليها باختزال شبيه بما سنصفه
أدناه، مع زيادة شرط: أن يحسب الاختزال في زمن كثير الحدود.

وقد جرى الاختزال ضمنًا فيما تقدم. فلنعرّف المجموعة~$K$ بوضع
\[
K = \Setabs{x}{\cfind{x}(x) \fdefined},
\]
أي $K=\Setabs{x}{x\in W_x}$. ويبين برهاننا على عدم قابلية مشكلة
التوقف للحل (\olref[hlt]{thm:halting-problem})، على نحو مباشر جدًا، أن
$K$ غير قابلة للحساب. وتذكر أن $K_0$ هي المجموعة
\[
K_0 = \Setabs{\tuple{e, x}}{\cfind{e}(x) \fdefined },
\]
أي $K_0=\Setabs{\tuple{e,x}}{x\in W_e}$. ومن السهل نقل البرهان على
عدم قابلية~$K$ للحساب إلى عدم قابلية~$K_0$ للحساب؛ إذ لو كانت
$K_0$ قابلة للحساب، لأمكننا تقرير ما إذا كان !!a{element}~$x$ في $K$
بالاكتفاء بالسؤال: هل الزوج $\tuple{x,x}$ في $K_0$. والدالة~$f$
التي تعين $x$ إلى $\tuple{x,x}$ مثال \emph{لاختزال} $K$ إلى $K_0$.
\end{explain}

\begin{defn}
لنأخذ مجموعتين من الأعداد الطبيعية $A$ و$B$. تسمى الدالة القابلة للحساب
~$f\colon\Nat\to\Nat$ \emph{اختزالًا متعددًا إلى واحد} لـ$A$ إلى~$B$
إذا تحقق، لكل عدد طبيعي~$x$، الشرط الآتي، ولا تسمى بذلك إلا حينئذ:
\[
x \in A \quad \text{إذا وفقط إذا} \quad f(x) \in B.
\]
فمتى وجد $f$ على هذه الصفة، كانت $A$ \emph{قابلة للاختزال متعددًا
إلى واحد} إلى~$B$، وكتبنا $A\leq_m B$. وإذا كانت $A$ قابلة للاختزال
متعددًا إلى واحد إلى $B$ والعكس، قيل إن $A$ و$B$ \emph{متكافئتان
متعددًا إلى واحد}، وكتبنا $A\equiv_m B$.
\end{defn}

فإن كانت $f$ في التعريف أحادية أيضًا، قلنا إن $A$
\emph{قابلة للاختزال واحدًا إلى واحد} إلى~$B$. وأكثر الاختزالات التي
سنوردها يفي بهذا الشرط الأقوى، وإن لم نحتج إليه في الاستدلال.

\begin{digress}
وقابلية الاختزال واحدًا إلى واحد أقوى حقًا
من قابلية الاختزال متعددًا إلى واحد، وإن لم يظهر ذلك لأول النظر. فثمة
مجموعتان لا نهائيتان $A$ و~$B$ تكون $A$ فيهما قابلة للاختزال متعددًا
إلى واحد إلى~$B$، ولكنها ليست قابلة للاختزال واحدًا إلى واحد إلى~$B$.
\end{digress}

\end{document}
